\chapter{Mackey completeness and Lipschitz curves}
\label{chap:mackey-lipschitz}

Glöckner and Neeb define Mackey completeness through Lipschitz curves (G--N~1.1.17(b)). In
this chapter we prove that their definition agrees with Definition~\ref{def:mackey-complete}.
The equivalence is stated in G--N~1.1.18(e) and \cite[2.14]{KM97}; the proof of the converse
below is ours.

\begin{definition}[Lipschitz curves]
  \label{def:lipschitz}
  \lean{LCAnalysis.diffQuot, LCAnalysis.IsLipschitzCurve}
  \leanok
  A curve $\gamma \colon [a, b] \to F$ is Lipschitz if the set of difference quotients
  \[
    \Bigl\{ \tfrac{1}{s - t}\bigl(\gamma(s) - \gamma(t)\bigr) : s, t \in [a, b],\ s \neq t \Bigr\}
  \]
  is bounded in $F$.
\end{definition}

\begin{definition}[Mackey completeness after Glöckner--Neeb]
  \label{def:integrates-lipschitz}
  \lean{LCAnalysis.IntegratesLipschitzCurves}
  \leanok
  \uses{def:lipschitz, def:weak-integral}
  We say that $F$ integrates Lipschitz curves if every Lipschitz curve $[a, b] \to F$, $a < b$,
  has a weak integral over $[a, b]$.
\end{definition}

\begin{lemma}[Lipschitz curves are continuous]
  \label{lem:lipschitz-continuous}
  \lean{LCAnalysis.IsLipschitzCurve.continuousOn, LCAnalysis.sub_mem_smul_convexHull_diffQuot}
  \leanok
  \uses{def:lipschitz}
  Let $\gamma$ be a Lipschitz curve on $[a, b]$ and $K = \conv(\{0\} \cup Q)$, where $Q$ is the
  set of difference quotients. Then $\gamma(s) - \gamma(t) \in h K$ whenever
  $0 \leq s - t \leq h$. In particular $\gamma$ is continuous.
\end{lemma}

\begin{proof}
  \leanok
  We have $\gamma(s) - \gamma(t) = (s - t)\, q$ with $q \in Q$, and $(s - t)q \in hK$ since $K$ is
  convex and contains $0$. Continuity follows because the bounded set $K$ is absorbed by every
  neighborhood of $0$.
\end{proof}

\section{Mackey complete spaces integrate Lipschitz curves}

\begin{theorem}
  \label{thm:mackey-to-lipschitz}
  \lean{LCAnalysis.MackeyComplete.exists_hasWeakIntegral_of_isLipschitzCurve, LCAnalysis.MackeyComplete.integratesLipschitzCurves}
  \leanok
  \uses{def:mackey-complete, def:integrates-lipschitz, lem:lipschitz-continuous, lem:riemann-limit, lem:conv-bounded}
  If $F$ is Mackey complete, then $F$ integrates Lipschitz curves.
\end{theorem}

\begin{proof}
  \leanok
  The set $K$ of Lemma~\ref{lem:lipschitz-continuous} is bounded by Lemma~\ref{lem:conv-bounded}.
  By Lemmas~\ref{lem:lipschitz-continuous} and~\ref{lem:riemann-refine},
  $R_{2^m}(\gamma) - R_{2^n}(\gamma) \in (b - a)\frac{b - a}{2^n}\, K$ for $n \leq m$, so the dyadic
  Riemann sums form a Mackey--Cauchy sequence. Its limit is the weak integral by
  Lemma~\ref{lem:riemann-limit}. This is the argument of Theorem~\ref{thm:exists-mackey}, with
  the difference quotients in place of the derivative.
\end{proof}

\section{The converse}

\begin{lemma}[Subsequences]
  \label{lem:mackey-subseq}
  \lean{LCAnalysis.tendsto_of_mackeyCauchy_of_subseq}
  \leanok
  \uses{def:mackey-complete}
  A Mackey--Cauchy sequence with a convergent subsequence converges, to the same limit.
\end{lemma}

\begin{proof}
  \leanok
  Let $u_{n_j} \to x$ and let $W$ be a neighborhood of $0$; choose a symmetric neighborhood $V$
  with $V + V \subseteq W$. For large $k$ we have $r_k B \subseteq V$. Choosing $j$ with
  $n_j \geq k$ and $u_{n_j} - x \in V$, we get
  $u_k - x = -(u_{n_j} - u_k) + (u_{n_j} - x) \in V + V \subseteq W$.
\end{proof}

The idea of the converse is to encode a Mackey--Cauchy sequence as a Lipschitz curve whose
integral is the limit we are looking for. The increments of the sequence are placed on disjoint
tents that accumulate at $t = 1$.

\begin{definition}[Tents]
  \label{def:tents}
  \lean{LCAnalysis.tentLeft, LCAnalysis.tentWidth, LCAnalysis.tent, LCAnalysis.tentSum, LCAnalysis.tentCurve}
  \leanok
  For $j \in \N$ let $l_j = 1 - 2^{-j}$ and $w_j = 2^{-(j+1)}$, so that the intervals
  $[l_j, l_j + w_j] = [l_j, l_{j+1}]$ tile $[0, 1)$. The $j$-th tent is the function
  \[
    \tau_j(t) = \max\Bigl(0,\ \frac{w_j}{2} - \Bigl|t - l_j - \frac{w_j}{2}\Bigr|\Bigr),
  \]
  supported on $[l_j, l_{j+1}]$, with slopes $\pm 1$ and height $w_j / 2$. For a sequence
  $(e_j)$ in $F$ we consider the curve $\gamma(t) = \sum_j \tau_j(t)\, e_j$, which is a finite
  sum for each $t$.
\end{definition}

\begin{lemma}[Properties of the tents]
  \label{lem:tents}
  \lean{LCAnalysis.integral_tent, LCAnalysis.abs_tent_sub_le, LCAnalysis.tent_nonneg, LCAnalysis.tent_le_one, LCAnalysis.continuous_tent, LCAnalysis.eventually_tentCurve_eq}
  \leanok
  \uses{def:tents}
  Each $\tau_j$ is continuous, takes values in $[0, 1]$, satisfies
  $|\tau_j(s) - \tau_j(t)| \leq |s - t|$, and $\int_0^1 \tau_j = w_j^2 / 4$. For each $t$,
  $\gamma(t) = \sum_{j < N} \tau_j(t)\, e_j$ for all large $N$.
\end{lemma}

\begin{proof}
  \leanok
  Elementary.
\end{proof}

\begin{theorem}
  \label{thm:lipschitz-to-mackey}
  \lean{LCAnalysis.MackeyComplete.of_integratesLipschitzCurves}
  \leanok
  \uses{def:integrates-lipschitz, def:mackey-complete, lem:tents, lem:mackey-subseq, lem:weak-integral-small, lem:weak-integral-algebra, lem:conv-bounded, lem:lipschitz-continuous}
  If $F$ integrates Lipschitz curves, then $F$ is Mackey complete.
\end{theorem}

\begin{proof}
  \leanok
  Let $(u_n)$ be Mackey--Cauchy, with $u_l - u_k \in r_k B$ for $k \leq l$. Let
  $K = \conv(\{0\} \cup B \cup -B)$, which is convex, symmetric and bounded. Put
  $c_j = 2^{-(j+1)}$, so that $\sum_j c_j = 1$.

  \emph{A fast subsequence.} Since $r_k \to 0$, we can choose $n_0 < n_1 < \cdots$ with
  $|r_{n_j}| \leq c_j w_j^2 / 4$. Let $d_j = u_{n_{j+1}} - u_{n_j}$ and
  $e_j = (4 / w_j^2)\, d_j$. Then $e_j \in c_j K$.

  \emph{The curve is Lipschitz.} Let $\gamma = \sum_j \tau_j e_j$. A difference quotient of
  $\gamma$ is a finite sum $\sum_j \frac{\tau_j(s) - \tau_j(t)}{s - t}\, e_j$ with coefficients
  of absolute value at most $1$, hence it lies in $\sum_j c_j K \subseteq K$. So $\gamma$ is a
  Lipschitz curve on $[0, 1]$, and by hypothesis it has a weak integral $y$.

  \emph{The integral is the limit.} Since $\int_0^1 \tau_j e_j = d_j$, the partial sum
  $\gamma_N = \sum_{j < N} \tau_j e_j$ has weak integral $u_{n_N} - u_{n_0}$. On the other hand,
  $\gamma(t) - \gamma_N(t) = \sum_{j \geq N} \tau_j(t) e_j \in 2^{-N} K$ for every $t$. By
  Lemma~\ref{lem:weak-integral-small}, applied with $W$ the closure of $2^{-N} K$, we get
  $y - (u_{n_N} - u_{n_0}) \in 2^{1-N}\, \overline{K}$, and this tends to $0$ since $K$ is
  bounded. Hence $u_{n_N} \to u_{n_0} + y$, and by
  Lemma~\ref{lem:mackey-subseq} the whole sequence converges.
\end{proof}

\begin{corollary}[The two definitions agree]
  \label{cor:mackey-iff}
  \lean{LCAnalysis.mackeyComplete_iff_integratesLipschitzCurves}
  \leanok
  \uses{thm:mackey-to-lipschitz, thm:lipschitz-to-mackey}
  A locally convex space is Mackey complete if and only if it integrates Lipschitz curves.
\end{corollary}

\begin{proof}
  \leanok
  Theorems~\ref{thm:mackey-to-lipschitz} and~\ref{thm:lipschitz-to-mackey}.
\end{proof}
